\chapter*{Введение}                         % Заголовок
\addcontentsline{toc}{chapter}{Введение}    % Добавляем его в оглавление

\newcommand{\actuality}{}
\newcommand{\progress}{}
\newcommand{\aim}{{\textbf\aimTXT}}
\newcommand{\tasks}{\textbf{\tasksTXT}}
\newcommand{\object}{{\textbf\objectTXT}}
\newcommand{\subject}{{\textbf\subjectTXT}}
\newcommand{\novelty}{\textbf{\noveltyTXT}}
\newcommand{\tinfluence}{\textbf{\tinfluenceTXT}}
\newcommand{\pinfluence}{\textbf{\pinfluenceTXT}}
\newcommand{\methods}{\textbf{\methodsTXT}}
\newcommand{\defpositions}{\textbf{\defpositionsTXT}}
\newcommand{\passport}{{\textbf\passportTXT}}
\newcommand{\reliability}{\textbf{\reliabilityTXT}}
\newcommand{\probation}{\textbf{\probationTXT}}
\newcommand{\contribution}{\textbf{\contributionTXT}}
\newcommand{\publications}{\textbf{\publicationsTXT}}

{\actuality} Задача декодирования состоит в восстановлении внешнего воздействия или его категории по отклику, зарегистрированному в виде многоканального сигнала высокой размерности. Формально требуется построить отображение из пространства сигналов в пространство воздействий по выборке, объем которой мал по сравнению с размерностью сигнала. В такой постановке определяющими становятся свойства сигнала, которые классические методы анализа многомерных временных рядов не используют.

Во-первых, отклик запаздывает относительно воздействия, и без явного учета лага сопоставление сигнала с воздействием смещено. Во-вторых, отклик сосредоточен в малой части каналов, поэтому информативные каналы необходимо отбирать до построения признаков. В-третьих, связи между каналами описываются ковариационной матрицей, которая принадлежит не линейному пространству, а многообразию симметричных положительно определенных матриц. Евклидова векторизация ковариаций игнорирует геометрию этого многообразия, а нейросетевые модели, способные выучить структуру сигнала из данных, требуют обучающих выборок на порядки больше доступных\extcite{shao2022transformers, boulmaiz2020impact}.

Риманова геометрия многообразия симметричных положительно определенных матриц с аффинно-инвариантной метрикой дает естественный инструмент для таких сигналов\extcite{moakher2005differential, arsigny2007geometric, bhatia2007positive}: расстояния инвариантны к преобразованию каналов, а проекция ковариаций на касательное пространство в точке риманова среднего переводит их в векторы, пригодные для линейных классификаторов\extcite{barachant2011multiclass, congedo2017riemannian, yger2016riemannian}. Однако проекция используется как эвристика: инвариантность признаков к преобразованию каналов, близость евклидовой геометрии касательного пространства к геометрии многообразия и устойчивость признаков к ошибке оценки ковариации по короткому ряду принимаются без явных констант. Получение таких гарантий необходимо для обоснованного выбора параметров метода и области его применимости.

Если сигнал регистрируется несколькими модальностями с разным разрешением по пространству и времени, лаг определяет и их синхронизацию, а объединение модальностей требует согласованных представлений. Совместный анализ дополняющих друг друга модальностей открывает возможности, недоступные каждой из них в отдельности\extcite{ritter2006simultaneous, huster2012methods}.

Многоканальные сигналы с такими свойствами возникают в разных областях: ковариационные дескрипторы изображений в компьютерном зрении\extcite{tuzel2008pedestrian, harandi2018dimensionality}, ковариации отраженного сигнала в радиолокации\extcite{arnaudon2013riemannian, formont2011statistical}, тензоры диффузии в медицинской визуализации\extcite{pennec2006riemannian}. В нейровизуализации эти особенности проявляются одновременно: функциональная магнитно-резонансная томография (фМРТ) дает высокое пространственное разрешение при низкой частоте сканирования и гемодинамическом лаге отклика\extcite{Glover2011, logothetis2008we}, электроэнцефалография (ЭЭГ)~--- высокое временное разрешение при слабой пространственной локализации\extcite{burle2015spatial}, а число размеченных объектов одного пациента измеряется десятками при размерности сигнала в десятки тысяч каналов\extcite{varoquaux2018cross}. Декодирование таких сигналов лежит в основе неинвазивных интерфейсов мозг--компьютер и применяется в диагностике\extcite{wolpaw2000brain, lebedev2017brain, nicolas2012brain, cassani2018systematic}. В работе методы строятся в общей постановке, а проверяются на данных нейровизуализации.

Таким образом, актуальна разработка пространственно-временных методов обработки сигнала высокой размерности, которые явно используют лаг отклика, локализацию информативных каналов и геометрию ковариационных связей между ними, согласуют модальности и сохраняют работоспособность при малом объеме выборки, а также строгое обоснование применяемых конструкций.

{\progress} Методы, использующие отдельные свойства многоканального сигнала, развивались независимо и преимущественно внутри приложений. Геометрия многообразия симметричных положительно определенных матриц с аффинно-инвариантной метрикой хорошо разработана\extcite{moakher2005differential, arsigny2007geometric, bhatia2007positive}: известны свойство метрического возрастания экспоненты, дающее нижнюю оценку геодезического расстояния через касательные векторы\extcite{bhatia2003exponential}, лог-евклидово приближение расстояния\extcite{arsigny2007geometric} и внутренняя граница Крамера--Рао для оценки ковариации\extcite{smith2005covariance}. Касательные признаки ковариаций как метод классификации многоканальных рядов сложились в компьютерном зрении\extcite{tuzel2008pedestrian} и интерфейсах мозг--компьютер\extcite{barachant2010riemannian, barachant2011multiclass, congedo2017riemannian, yger2016riemannian}; явные оценки искажения расстояний проекцией и неасимптотические оценки шума признаков там не приводятся, а к сигналам с десятками тысяч каналов и выборкой в десятки объектов метод не применялся: в анализе фМРТ риманова геометрия служит лишь описанию связности\extcite{you2021re}.

Лаг отклика в моделях, связывающих воздействие и сигнал, задается фиксированным; для гемодинамического отклика изучены его природа\extcite{Bandettini1992, Logothetis2003} и изменчивость между источниками и областями\extcite{handwerker2004variation, handwerker2012continuing}, но влияние выбора лага на качество моделей отклика исследовано мало. Отбор информативных каналов по связи с воздействием восходит к анализу областей интереса\extcite{poldrack2007region, haxby2001distributed} и выбору каналов по римановым расстояниям\extcite{barachant2011channel}, однако выполняется до классификации и единообразно для всех классов.

Реконструкция воздействия по сигналу высокой размерности опирается на контрастивное обучение представлений\extcite{radford2021learning, wang2023connecting} и диффузионные генеративные модели\extcite{ho2020denoisingdiffusionprobabilisticmodels}; для одной модальности такие методы доведены до реконструкции изображений и видео\extcite{ozcelik2023naturalscenereconstructionfmri, scotti2023reconstructingmindseyefmritoimage, chen2023cinematicmindscapeshighqualityvideo, li2024visual} и требуют больших обучающих выборок, а объединение двух синхронно регистрируемых модальностей в одной модели реконструкции не рассматривалось, поскольку открытые наборы одновременных записей появились недавно\extcite{Telesford2023}. Восстановление пропущенных каналов выполняется интерполяцией по геометрии сенсоров\extcite{perrin1989spherical} либо общими вероятностными моделями пропусков во временных рядах\extcite{tashiro2021csdi}, не использующими расположение сенсоров.

{\aim} данной работы является разработка пространственно-временных методов обработки сигналов высокой размерности, обеспечивающих декодирование внешнего воздействия в условиях ограниченного объема выборки. Для достижения цели были поставлены и решены следующие {\tasks}:
\begin{enumerate}[beginpenalty=10000]
    \item Построить модель прогноза отклика по воздействию, обратную к декодированию, и на ее основе метод оценки лага отклика.
    \item Разработать метод классификации многоканальных временных рядов в условиях ограниченной выборки на основе масок активности классов и проекции ковариационных матриц на касательное пространство риманова многообразия.
    \item Исследовать теоретические свойства касательных признаков ковариационных матриц: инвариантность к преобразованию каналов, искажение расстояний при проекции и устойчивость к шуму оценки ковариации.
    \item Построить мультимодальную архитектуру реконструкции воздействия по двум синхронно регистрируемым сигналам и метод восстановления отсутствующих каналов многоканального сигнала по геометрии сенсоров.
\end{enumerate}

{\object} настоящей работы являются многоканальные сигналы высокой размерности, регистрируемые в ответ на внешнее воздействие.

{\subject} данной работы являются методы обработки таких сигналов, обеспечивающие декодирование внешнего воздействия: восстановление пропущенных каналов, оценка лага отклика, отбор каналов и построение ковариационных признаков, согласование модальностей, а также теоретические свойства этих методов.

{\novelty}
\begin{enumerate}[beginpenalty=10000]
    \item Предложен метод оценки лага отклика по минимуму ошибки авторегрессионного прогноза отклика по воздействию; показано, что в локализованной области отклика ошибка имеет выраженный минимум, согласованный с независимыми оценками лага, тогда как по полному сигналу минимум не выражен.
    \item Предложен метод классификации многоканальных временных рядов, в котором объект описывается отдельно с точки зрения каждого класса: маска активности класса отбирает каналы по кросс-корреляции с рядом воздействия, а ковариация их рядов проецируется на касательное пространство в точке риманова среднего; размерность признаков не зависит от объема выборки.
    \item Доказаны инвариантность касательных признаков к невырожденному преобразованию каналов, двусторонняя оценка искажения расстояний с неулучшаемой константой, выраженной через спектральный разброс касательных векторов, и оценка концентрации ошибки оценки ковариации в геодезической метрике, распределение которой не зависит от истинной ковариации.
    \item Предложена мультимодальная архитектура реконструкции воздействия по двум синхронно регистрируемым сигналам с контрастивным выравниванием совместных представлений с пространством предобученных эмбеддингов воздействия и двухстадийной диффузионной генерацией; показано, что объединение модальностей повышает качество реконструкции по сравнению с каждой из них в отдельности. Для записей с неполным набором каналов предложен метод восстановления отсутствующих каналов по геометрии расположения сенсоров и его вероятностное обобщение на основе маскированного согласования потоков с доказанной корректностью.
\end{enumerate}

{\tinfluence} полученных результатов заключается в строгом обосновании широко применяемой конструкции касательных признаков ковариационных матриц: доказаны инвариантность признаков к преобразованию каналов, двусторонняя оценка искажения расстояний с явной константой и оценка шума оценки ковариации, связывающая число каналов и длину ряда. Результаты не зависят от природы сигнала и применимы к любым многоканальным временным рядам, описываемым ковариацией каналов.

{\pinfluence} настоящей работы заключается в применимости методов к произвольным многоканальным сигналам высокой размерности и, в частности, методов оценки лага и классификации к выборке одного источника объемом в десятки объектов, где нейросетевые модели неприменимы, и в нейросетевой архитектуре реконструкции воздействия, обучаемой на синхронных записях многих источников; методы проверены на открытых наборах данных по воспроизводимым протоколам, а их программные реализации опубликованы в открытых репозиториях. Результаты анализа лага учтены при построении масок активности и синхронизации модальностей, а метод восстановления отсутствующих каналов позволяет применять обученные модели к записям с неполным набором сенсоров.

{\methods} Для достижения поставленных в диссертационной работе целей используются:
\begin{enumerate}[beginpenalty=10000]
    \item Методы линейной алгебры и римановой геометрии: аффинно-инвариантная метрика на многообразии симметричных положительно определенных матриц, матричные функции, дифференциал матричной экспоненты.
    \item Методы теории вероятностей и математической статистики: неасимптотические оценки сингулярных чисел гауссовых матриц, проверка статистических гипотез с поправкой на множественные сравнения, доверительные интервалы.
    \item Методы машинного обучения и вычислительного эксперимента: регуляризованная линейная регрессия, логистическая регрессия, контрастивное обучение представлений, диффузионные генеративные модели, воспроизводимая постановка экспериментов на открытых данных.
\end{enumerate}

{\defpositions}
\begin{enumerate}[beginpenalty=10000]
    \item Линейная авторегрессионная модель прогнозирования отклика на внешнее воздействие и метод оценки лага отклика на ее основе.
    \item Метод классификации многоканальных временных рядов на основе масок активности классов и проекции ковариационных матриц на касательное пространство риманова многообразия, устойчивый в условиях ограниченного объема выборки.
    \item Теоретические свойства проекции ковариационных матриц на касательное пространство в точке риманова среднего: инвариантность признаков к невырожденному преобразованию каналов, двусторонняя оценка искажения расстояний с неулучшаемой константой и оценка ошибки оценки ковариации по длине ряда.
    \item Мультимодальная архитектура реконструкции воздействия по двум синхронно регистрируемым сигналам с контрастивным выравниванием совместных представлений и двухстадийной диффузионной генерацией.
\end{enumerate}

{\passport} Диссертационная работа соответствует паспорту научной специальности 1.2.1~--- <<Искусственный интеллект и машинное обучение>>, а именно в наибольшей степени следующим направлениям исследований: п.~4 <<Разработка методов, алгоритмов и создание систем искусственного интеллекта и машинного обучения для обработки и анализа текстов на естественном языке, для изображений, речи, биомедицины и других специальных видов данных>>; п.~12 <<Исследования в области надежности, устойчивости и переобучения систем класса ИИ>>; п.~15 <<Математические исследования в области статистики, логики, алгебры, топологии, анализа функций и других областях, ориентированные на решение задач искусственного интеллекта и машинного обучения>>; п.~16 <<Исследования в области специальных методов оптимизации, проблем сложности и снижения размерности>>; п.~17 <<Исследования в области многослойных алгоритмических конструкций, в том числе многослойных нейросетей>>.

{\reliability} полученных результатов обеспечивается корректностью применения апробированного математического аппарата римановой геометрии, теории вероятностей и математической статистики, математически строгими доказательствами теоретических утверждений, а также экспериментальной проверкой разработанных методов на открытых наборах данных с фиксированными протоколами разбиения, статистическими критериями с поправкой на множественные сравнения и доверительными интервалами. Для вычислительных экспериментов даются детальные описания и опубликован программный код, обеспечивающие их воспроизводимость. Результаты опубликованы в рецензируемых научных изданиях и доложены на российских и международных конференциях по машинному обучению и анализу данных.

{\probation} Основные результаты работы докладывались и обсуждались на следующих научных конференциях:
\begin{enumerate}
    \item 16-я Международная конференция <<Интеллектуализация обработки информации (ИОИ-2026)>>,
    Витебск, Беларусь, 28 сентября -- 1 октября 2026~г.
    \item 68-я Всероссийская научная конференция МФТИ,
    Долгопрудный, 30 марта -- 4 апреля 2026~г.
    \item X Международная конференция по искусственному интеллекту AI Journey 2025,
    Москва, 19--21 ноября 2025~г.
    \item 22-я Всероссийская конференция с международным участием
    <<Математические методы распознавания образов (ММРО-2025)>>,
    Муром, 22--26 сентября 2025~г.
    \item 67-я Всероссийская научная конференция МФТИ,
    Долгопрудный, 31 марта -- 5 апреля 2025~г.
    \item XXVI Международная научно-техническая конференция <<Нейроинформатика-2024>>,
    Долгопрудный, 21--25 октября 2024~г.
    \item 66-я Всероссийская научная конференция МФТИ,
    Долгопрудный, 1--6 апреля 2024~г.
\end{enumerate}


{\contribution} Все приведенные результаты, кроме отдельно оговоренных случаев, получены диссертантом лично при научном руководстве к.ф.-м.н.~А.\,В.~Грабового.

\ifnumequal{\value{bibliosel}}{0}
{%%% Встроенная реализация с загрузкой файла через движок bibtex8. (При желании, внутри можно использовать обычные ссылки, наподобие `\cite{vakbib1,vakbib2}`).
    {\publications} Основные результаты по теме диссертации изложены
    в~XX~печатных изданиях,
    X из которых изданы в журналах, рекомендованных ВАК,
    X "--- в тезисах докладов.
}%
{%%% Реализация пакетом biblatex через движок biber
    \begin{refsection}[bl-author, bl-registered]
        % Это refsection=1.
        % Процитированные здесь работы:
        %  * подсчитываются, для автоматического составления фразы "Основные результаты ..."
        %  * попадают в авторскую библиографию, при usefootcite==0 и стиле `\insertbiblioauthor` или `\insertbiblioauthorgrouped`
        %  * нумеруются там в зависимости от порядка команд `\printbibliography` в этом разделе.
        %  * при использовании `\insertbiblioauthorgrouped`, порядок команд `\printbibliography` в нем должен быть тем же (см. biblio/biblatex.tex)
        %
        % Невидимый библиографический список для подсчета количества публикаций:
        \phantom{\printbibliography[heading=nobibheading, section=1, env=countauthorvak,          keyword=biblioauthorvak]%
        \printbibliography[heading=nobibheading, section=1, env=countauthorwos,          keyword=biblioauthorwos]%
        \printbibliography[heading=nobibheading, section=1, env=countauthorscopus,       keyword=biblioauthorscopus]%
        \printbibliography[heading=nobibheading, section=1, env=countauthorconf,         keyword=biblioauthorconf]%
        \printbibliography[heading=nobibheading, section=1, env=countauthorother,        keyword=biblioauthorother]%
        \printbibliography[heading=nobibheading, section=1, env=countauthor,             keyword=biblioauthor]%
        \printbibliography[heading=nobibheading, section=1, env=countauthorvakscopuswos, filter=vakscopuswos]%
        \printbibliography[heading=nobibheading, section=1, env=countauthorscopuswos,    filter=scopuswos]}%
        %
        \nocite{*}%
        %
        {\publications} \ifsynopsis\textit{Список публикаций приведен в конце автореферата. }\else\textit{Публикации автора включены в общий список литературы. }\fi Основные результаты по теме диссертации изложены в~\arabic{citeauthor}~научных работах, из которых~\arabic{citeauthorscopuswos} "--- в~периодических научных журналах, индексируемых Web of~Science и Scopus, а остальные~\arabic{citeauthorconf} "--- в~тезисах докладов российских и международных конференций.
    \end{refsection}%
    \begin{refsection}[bl-author, bl-registered]
        % Это refsection=2.
        % Процитированные здесь работы:
        %  * попадают в авторскую библиографию, при usefootcite==0 и стиле `\insertbiblioauthorimportant`.
        %  * ни на что не влияют в противном случае
        \nocite{dorin2024forecasting}
        \nocite{dorin2025enhancing}
        \nocite{dorin2026decoding}
        \nocite{dorin2025mmro}
        \nocite{dorin2024conf66mipt}
        \nocite{dorin2025conf67mipt}
        \nocite{dorin2026conf68mipt}
        \nocite{dorin2025aij}
        \nocite{dorin2026ioi}
    \end{refsection}%
}
 % Характеристика работы по структуре во введении и в автореферате не отличается (ГОСТ Р 7.0.11, пункты 5.3.1 и 9.2.1), потому её загружаем из одного и того же внешнего файла, предварительно задав форму выделения некоторым параметрам

\textbf{Объем и структура работы.} Диссертация состоит из~введения,
\formbytotal{totalchapter}{глав}{ы}{}{} и заключения.
%% на случай ошибок оставляю исходный кусок на месте, закомментированным
%Полный объем диссертации составляет  \ref*{TotPages}~страницу
%с~\totalfigures{}~рисунками и~\totaltables{}~таблицами. Список литературы
%содержит \total{citenum}~наименований.
%
Полный объем диссертации составляет
\formbytotal{TotPages}{страниц}{у}{ы}{}, включая
\formbytotal{totalcount@figure}{рисун}{ок}{ка}{ков} и
\formbytotal{totalcount@table}{таблиц}{у}{ы}{}.
Список литературы содержит \formbytotal{citenum}{наименован}{ие}{ия}{ий}.

В главе~1 дана общая постановка задачи декодирования, описаны пространственно-временные характеристики многоканальных сигналов и методы декодирования по одной модальности, рассмотрена область применения. В главе~2 строится модель прогноза сигнала по стимулу, обратная к декодированию, и на ее основе оценивается лаг сигнала. В главе~3 предлагается метод классификации многоканальных временных рядов по малой выборке на основе масок активности и проекции ковариаций на касательное пространство и доказываются свойства касательных признаков. В главе~4 строится мультимодальная архитектура реконструкции стимула по двум синхронно регистрируемым сигналам и метод восстановления отсутствующих каналов. В заключении сформулированы основные результаты работы.
% \formbytotal{citenum}{наименован}{ие}{ия}{ий}.